% -----------------------------------------------------------------------------
% This file is automatically generated.
% Do not edit manually.
% Source: notebooks/support/regression-analysis/support.Rmd
% -----------------------------------------------------------------------------


\noindent We can assess normality issues with a histogram or with a QQ plot. The histogram in Figure 5.14 is obtained by:

\par\addvspace{\topsep}
\begingroup
\fvset{listparameters={%
  \setlength{\topsep}{0pt}%
  \setlength{\partopsep}{0pt}%
  \setlength{\parsep}{0pt}%
  \setlength{\itemsep}{0pt}%
}}
\begin{Verbatim}[breaklines=true,formatcom=\color{ada_blue}]
# Calculate normalized residuals
nres <- ussteamco.mod.3$residuals / sd(ussteamco.mod.3$residuals)
model3 <- as.data.frame(nres)
hist_stand_res <-
  ggplot(model3, aes(nres)) +
  geom_histogram(aes(y = after_stat(density)),
    breaks = seq(-2.55, 1.65, by = 0.7),
    fill = "#00338D"
  ) +
  stat_function(fun = dnorm, args = list(
    mean = mean(model3$nres),
    sd = sd(model3$nres)
  ))
hist_stand_res
\end{Verbatim}
\endgroup
\par\addvspace{\topsep}

\noindent The QQ plot is produced by:
